
\chapter{Cистемийн онолын судалгаа}
\label{Chapter1}

\section{Цахим сургалтын системийн тухай ойлголт}
Цахим сургалтын систем нь суралцах үйл ажиллагааг программ хангамжийн орчинд төлөвлөх, түгээх, хянах, үнэлэх, сайжруулахад чиглэсэн мэдээллийн систем юм. Ийм систем нь хичээлийн агуулга, үнэлгээ, хэрэглэгчийн профайл, суралцах ахиц, харилцаа холбоо, тайлан шинжилгээ зэрэг олон давхар мэдээллийг нэгтгэн удирддаг. Орчин үеийн e-learning платформууд нь зөвхөн файлууд хадгалах эсвэл бичлэг байрлуулах хэрэгсэл байхаа больж, суралцагчийн суралцах процессын идэвхтэй хэсэг болж өөрчлөгдөж байна \cite{elearning}.

Цахим сургалтыг уламжлалт танхимын сургалтаас ялгах хэд хэдэн онцлог бий.
\begin{itemize}
\item Агуулгыг байршил, цаг хугацаанаас үл хамааран түгээх боломжтой.
\item Суралцагч өөрийн хурдаар давтаж үзэх, буцаан харах боломжтой.
\item Хэрэглэгчийн үйлдэл бүрийг бүртгэж, өгөгдөлд суурилсан дүн шинжилгээ хийх боломжтой.
\item Автомат үнэлгээ, санал болголт, ухаалаг туслах зэрэг нэмэлт үйлчилгээ нэгтгэхэд хялбар.
\end{itemize}

Гэсэн хэдий ч цахим сургалтын энгийн системүүдэд дараах сул талууд нийтлэг ажиглагддаг.
\begin{itemize}
\item Бүх хэрэглэгчид ижил агуулга, ижил дарааллаар хүрдэг.
\item Хэрэглэгчийн ойлголтын түвшин, өмнөх мэдлэг, алдааны хэв маягийг тооцохгүй.
\item Багш шууд оролцоогүй үед тайлбар, чиглэл, тусламж авах боломж сул.
\item Хичээлээ дундаас орхих, идэвх буурах, ахиц зогсох эрсдэл өндөр.
\end{itemize}

Иймээс орчин үеийн сургалтын системийн хөгжлийн чиглэл нь хувь хүнд тохирсон агуулга, өгөгдөлд суурилсан ахицын шинжилгээ, автоматжуулсан тайлбар, шалгалт, санал болголтыг нэгтгэхэд чиглэж байна.

\section{Програмчлалын сургалтын онцлог}
Програмчлалын сургалт нь бусад онолын хичээлээс хэд хэдэн зүйлээр ялгаатай. Нэгдүгээрт, мэдлэг нь хуримтлагдах шинжтэй. Өмнөх ойлголт сул байвал дараагийн агуулгыг ойлгох боломж эрс буурдаг. Хоёрдугаарт, зөвхөн уншиж ойлгох бус, бодлого бодох, код бичих, алдааг засах, гүйцэтгэлийг шалгах зэрэг практик үйл ажиллагаа зайлшгүй шаардлагатай. Гуравдугаарт, суралцагчдын анхны түвшин хоорондын ялгаа их байдаг.

Програмчлалын сургалтын системд дараах шаардлага онцгой ач холбогдолтой.
\begin{itemize}
\item Суурь ойлголтыг баталгаажуулах шаталсан агуулгын бүтэц
\item Богино хугацаанд ойлголтоо шалгах шалгалт болон дасгал
\item Алдаа гарсан үед тайлбар болон зөвлөмж өгөх туслах орчин
\item Суралцагчийн ахицыг тоон үзүүлэлтээр харах боломж
\item Онол, практик, харилцаа, урамшууллыг нэгтгэсэн хэрэглэгчийн туршлага
\end{itemize}

Эдгээр хэрэгцээг нэг системд цогцоор нь шийдэхийн тулд Дасан зохицох сургалт болон Хиймэл оюуны тусламжтай сургалт-ийн аргачлалуудыг нэгтгэх шаардлага үүсдэг.

\section{Дасан зохицох сургалт-ийн үндэс}
Дасан зохицох сургалт нь суралцагч бүрийн онцлогт тохируулан агуулга, даалгавар, үнэлгээ, дараагийн алхмыг өөрчилдөг арга зүй юм. Энэ нь статик curriculum-ээс ялгаатайгаар өгөгдөлд тулгуурласан шийдвэр гаргалтыг дэмждэг. Хэрэглэгчийн Түвшин тогтоох шалгалт-ийн үр дүн, Хичээл гүйцэтгэл, Тестийн оноо, нийт оноо, агуулга ашигласан давтамж зэрэг мэдээлэл нь дараагийн санал болголтыг тодорхойлох оролт болдог.

Дасан зохицох сургалт системд түгээмэл хэрэглэгддэг үндсэн механизмууд:
\begin{enumerate}
\item Анхны түвшин тогтоох механизм
\item Агуулгыг түвшин, сэдвээр ангилах бүтэц
\item Явцын өгөгдөл цуглуулах логик
\item Нөхцөлд суурилсан дараагийн алхмын шийдвэр
\item Ахиц, амжилтын үзүүлэлтээр дахин үнэлэх процесс
\end{enumerate}

Хэрэв хэрэглэгч анхан түвшинд байвал шууд ахисан агуулга санал болгох нь тохиромжгүй. Харин шат дараалсан байдлаар суурь мэдлэгийг бататгаж, бага багаар хүндрэл нэмэгдүүлэх нь илүү үр дүнтэй. Үүний эсрэгээр ахисан хэрэглэгчийг хэт энгийн агуулгаар удаан саатуулах нь суралцах идэвхийг бууруулна. Дасан зохицох сургалт нь энэ тэнцвэрийг хадгалах зорилготой.

\subsection{Хувь хүнд тохирсон зөвлөмж}
зөвлөмж логик нь Дасан зохицох сургалт-ийн практик илэрхийлэл болдог. Хэрэглэгчийн түвшин, дуусгасан хичээл, оноо, сонирхсон контент, like дарсан хичээл, орсон track зэрэг мэдээлэлд тулгуурлан дараагийн курс эсвэл хичээл-ийг санал болгоно. Энгийн нөхцөлд rule-based зөвлөмж ашиглаж болох бөгөөд илүү өргөн өгөгдөлтэй үед machine сургалт дээр суурилсан зөвлөмж руу шилжих боломжтой.

Энэхүү дипломын төслийн хүрээнд зөвлөмж нь дараах зарчимд тулгуурласан.
\begin{itemize}
\item Хэрэглэгчийн түвшинтэй тэнцүү эсвэл түүнээс доод түвшний course-уудыг харуулах
\item Дуусгасан хичээл-үүдийг санал болгосон жагсаалтаас хасах
\item Хэрэглэгчийн одоогийн түвшинд тохирсон агуулгыг тэргүүн ээлжид санал болгох
\item Курс болон суралцах замналыг эрэмбэлэхдээ хэрэглэгчийн идэвх, хичээлийн бүтэц, агуулгын дарааллыг харгалзах
\end{itemize}

\section{Үүсгэгч хиймэл оюун-ийн сургалтын хэрэглээ}
Үүсгэгч хиймэл оюун нь текст, зураг, код, асуулт, тайлбар зэрэг шинэ контент үүсгэх чадвартай бөгөөд сургалтын системд дараах үндсэн хэрэглээтэй.
\begin{itemize}
\item Хичээлийн урт агуулгыг богино, ойлгомжтой хураангуй болгон хувиргах
\item Ойлголт шалгах асуулт, сонголттой тест боловсруулах
\item Практик даалгавар, жишээ бодлого, тайлбар бичих
\item Чатбот хэлбэрээр асуулт хариулт хийх
\item Хэрэглэгчийн оруулсан агуулгыг moderation хийх
\end{itemize}

Энэ төрлийн Хиймэл оюуны тусламж нь багшийн оролцоог бүрэн орлохгүй ч давтагддаг, их цаг шаарддаг, шуурхай байх ёстой үйл ажиллагааг дэмжиж чадна. Жишээлбэл нэг хичээл-ийн агуулгаас 1--2 өгүүлбэрийн summary, 3 асуулттай quiz, эсвэл нэг жижиг practice task автоматаар үүсгэж өгөх нь хэрэглэгчийн туршлагыг мэдэгдэхүйц сайжруулна. Мөн хэрэглэгч ойлгоогүй хэсгээ асуухад чатбот тайлбар өгөх боломжтой болдог.

\subsection{Хиймэл оюун ашиглахад анхаарах асуудлууд}
Хиймэл оюун нэгтгэхдээ зөвхөн боломжид төвлөрөх бус, эрсдэл, хязгаарлалтыг мөн тооцох шаардлагатай.
\begin{itemize}
\item Үүсгэсэн агуулга буруу эсвэл тодорхой бус байж болно.
\item Гадаад Хиймэл оюуны үйлчилгээ тасалдах, хурд удаашрах, API key асуудал гарах боломжтой.
\item Хэрэглэгчийн өгөгдөл, нууцлал, зохисгүй контентын эрсдэл бий.
\item Хэт их Хиймэл оюуны хамааралтай бол системийн үндсэн функц тасалдах магадлалтай.
\end{itemize}

Иймээс Хиймэл оюуны тусламжтай сургалт систем нь fallback логик, error handling, агуулгын хязгаарлалт, moderation механизмтай байх ёстой. Энэхүү дипломын төсөлд summary болон practice task хэсэгт fallback арга ашигласнаар Хиймэл оюуны үйлчилгээ доголдсон үед ч системийн гол урсгал бүрэн зогсохгүй байхаар төлөвлөсөн.

\section{Веб архитектур ба full-stack хөгжүүлэлт}

Орчин үеийн сургалтын платформыг хөгжүүлэхэд frontend, backend, өгөгдлийн сан, authentication, файлын менежмент, API интеграц, байршуулалт зэрэг олон давхар бүрэлдэхүүн шаардлагатай.

Frontend нь хэрэглэгчийн интерфейс, маршрутын удирдлага, локал төлөв  удирдах, формын шалгалт, анимаци болон API дуудлагыг хариуцна.

Backend нь бизнес логик, өгөгдлийн баталгаажуулалт, эрхийн түвшин, өгөгдөл сериализаци, authentication, хиймэл оюун холбогдох endpoint-ууд, хэрэглэгчийн ахиц дүн, контент хяналт зэргийг удирдана.

Frontend технологийн хувьд React нь component-based бүтэц, өргөтгөх чадвар, олон хуудсатай интерфейсийг илүү зохион байгуулалттай хөгжүүлэх боломж олгодог \cite{react}. TypeScript нь өгөгдлийн төрөл тодорхой байлгаж, API харилцааны алдааг эрт илрүүлэх, томрох боломжтой кодын санг илүү тогтвортой байлгахад тустай \cite{typescript}. Vite нь хурдан хөгжүүлэлтийн сервер, орчин үеийн build pipeline, frontend developer experience-ийг сайжруулдаг \cite{vite}.

Backend-ийн хувьд Django болон Django REST Framework нь model-view-serializer бүтэц, authentication, admin, ORM, REST API гаргах боломжийг нэг дор өгдөг \cite{django,drf}. Ийм төрлийн системд өгөгдлийн бүтэц олон холбоостой байдаг тул ORM ашиглах нь хөгжүүлэлтийн хурд, өргөтгөх чадварыг сайжруулна. JWT authentication нь frontend болон backend тусдаа байрлаж буй архитектурт хэрэглэгчийн нэвтрэлтийг удирдахад түгээмэл хэрэглэгддэг \cite{jwt}.

\section{Өгөгдөл хадгалалт ба нэвтрүүлэлтийн орчин}
Сургалтын систем нь хэрэглэгчийн бүртгэл, курс, хичээл, ахиц дүн, community контент, Хиймэл оюун-аар үүсгэсэн туслах контент зэрэг олон төрлийн өгөгдөл хадгалдаг тул өгөгдлийн сангийн тогтвортой байдал чухал. PostgreSQL нь харилцаат өгөгдлийн сангийн найдвартай, өргөн хэрэглээтэй систем бөгөөд production орчинд ашиглахад тохиромжтой \cite{postgres}. Харин хөгжүүлэлтийн үед SQLite ашиглах боломжтой байхаар орчны тохиргоо хийх нь хөгжүүлэгчдийн ажлыг хөнгөвчилнө.

Нэвтрүүлэлт болон хөгжүүлэлтийн орчны ялгааг багасгахын тулд Docker ашиглах нь үр дүнтэй \cite{docker}. Docker Compose ашигласнаар backend болон database service-үүдийг нэг файл дээр тодорхойлж, орчны хувьсагч, сүлжээ, эзлэхүүн, health check-ийг удирдах боломжтой. Энэ нь ``миний компьютер дээр ажиллаж байна'' гэсэн түгээмэл асуудлыг багасгаж, системийг багийн түвшинд тогтвортой ажиллуулахад ач холбогдолтой.

\section{Сургалтын систем дэх community ба тоглоомжуулалт}
Суралцах процесс нь зөвхөн материал унших, шалгалт өгөхөөс бүрдэхгүй. Хэрэглэгчид хоорондоо асуулт асуух, туршлага солилцох, бие биеэ дэмжих, сонирхлыг тогтвортой хадгалах хэрэгцээтэй байдаг. Community module нь хичээл-тэй холбоотой хэлэлцүүлэг, нийтлэл, сэтгэгдлийн урсгалыг дэмждэг. Энэ нь хэрэглэгч системтэй зөвхөн ``контент хэрэглэгч'' байдлаар бус ``оролцогч'' байдлаар харилцах боломж олгоно.

Gamification буюу тоглоомжуулалт нь урамшуулал, ахицын мэдрэмж, жижиг сорилт, зорилт, сонирхол татах элемент ашиглан суралцах үйл явцыг идэвхжүүлдэг. Энэхүү дипломын системд games хэсгийг оруулснаар хэрэглэгч хичээлийн урсгалтай холбоотой интерактив контент ашиглах боломжтой болсон. Энэ нь ялангуяа урт хугацаанд системийг идэвхтэй ашиглах зан төлөвт эерэг нөлөө үзүүлнэ.

\section{Судлагдсан байдлын нэгтгэл}
Судалгааны үр дүнгээс дараах дүгнэлтүүд гарч байна.
\begin{itemize}
\item Цахим сургалтын системийн үнэ цэнэ нь зөвхөн агуулгын хадгалалт бус, агуулгыг ухаалаг удирдах чадвараар хэмжигдэх болсон.
\item Програмчлалын сургалт нь түвшин тогтоох, дасгал ажиллуулах, алдааг тайлбарлах, шаталсан агуулга санал болгох механизмыг онцгой шаардана.
\item Дасан зохицох сургалт болон Хувь хүнд тохирсон зөвлөмж нь хэрэглэгчийн түвшин, ахицад тулгуурласан бодит шаардлага юм.
\item Үүсгэгч хиймэл оюун нь хураангуй, асуулт, дадлагын даалгавар, чатбот, хяналт зэрэг олон үйлдэлд ашиглагдах боломжтой.
\item Fallback логик, эрхийн удирдлага, өгөгдлийн бүтэц, байршуулалт орчин зэрэг инженерчлэлийн асуудлуудыг хамтад нь шийдэхгүйгээр ийм системийг тогтвортой хэрэгжүүлэх боломжгүй.
\end{itemize}

\section{Бүлгийн дүгнэлт}
Энэ бүлэгт цахим сургалт, програмчлалын сургалтын онцлог, Дасан зохицох сургалт, Хиймэл оюуны тусламжтай сургалт, full-stack веб архитектур, өгөгдлийн сан, байршуулалт орчин, community болон gamification-ийн онолын үндсийг авч үзлээ. Судалгаанаас харахад хэрэглэгчийн түвшин, ахиц, идэвхэд суурилсан сургалтын систем нь өнөөгийн хэрэгцээнд нийцсэн бөгөөд Хиймэл оюуны интеграцитай хослуулах нь сургалтын үр өгөөж, хүртээмж, хэрэглэгчийн туршлагыг мэдэгдэхүйц сайжруулах боломжтой. Энэ онолын үндэслэл нь дараагийн бүлэгт авч үзэх системийн шинжилгээ, шаардлага тодорхойлолтын суурь болж байна.
